% ch26: 数项级数
\chapter{数项级数——无穷项相加}
\begin{applicationbox}
学完本章：理解级数收敛与发散、正项级数判别法
\end{applicationbox}
\section{级数的概念} \(\sum_{n=1}^\infty a_n = \lim_{n\to\infty} S_n\)，\(S_n=a_1+\cdots+a_n\)。

收敛当 \(\lim S_n\) 存在，用 ε-N 语言：对任意 \(\varepsilon>0\)，存在 \(N\)，
使得 \(m>n>N\) 时 \(|S_m - S_n| = |a_{n+1}+\cdots+a_m| < \varepsilon\)（柯西准则）。
\section{正项级数判别法} 比较判别法、比值判别法、根值判别法、积分判别法。

\textbf{比较判别法的 ε-N 证明：} 设 \(0\le a_n\le b_n\)，\(\sum b_n\) 收敛于 \(S\)。
部分和 \(B_n=b_1+\cdots+b_n\)，由 ε-N 定义，对任意 \(\varepsilon>0\)，存在 \(N\)，
使 \(m>n>N\) 时 \(|B_m-B_n|=b_{n+1}+\cdots+b_m<\varepsilon\)。
由于 \(a_n\le b_n\)，有 \(A_m-A_n=a_{n+1}+\cdots+a_m\le b_{n+1}+\cdots+b_m<\varepsilon\)，
故 \(\{A_n\}\) 是柯西列，\(\sum a_n\) 收敛。\(\blacksquare\)
\begin{examplebox}[1] 几何级数 \(\sum r^n\)：\(|r|<1\) 收敛，\(|r|\ge1\) 发散。
\(p\)-级数 \(\sum 1/n^p\)：\(p>1\) 收敛，\(p\le1\) 发散。\end{examplebox}
\section{习题}
\begin{exercise}[0] 级数收敛的定义是什么？\end{exercise}
\begin{exercise}[0] 几何级数 \(\sum r^n\) 收敛的条件？\end{exercise}
\begin{exercise}[0] \(p\)-级数 \(\sum 1/n^p\) 收敛的条件？\end{exercise}
\begin{exercise}[0] 若通项不趋于0，级数一定发散吗？\end{exercise}
\begin{exercise}[1] 判断 \(\sum \frac1{2^n}\) 的收敛性。\end{exercise}
\begin{exercise}[1] 判断 \(\sum \frac1{n}\) 的收敛性。\end{exercise}
\begin{exercise}[1] 判断 \(\sum \frac1{n^2}\) 的收敛性。\end{exercise}
\begin{exercise}[1] 求 \(\sum_{n=0}^\infty (\frac12)^n\) 的和。\end{exercise}
\begin{exercise}[2] 用比值法判断 \(\sum \frac{n}{2^n}\)。\end{exercise}
\begin{exercise}[2] 用比较法判断 \(\sum \frac1{n^2+1}\)。\end{exercise}
\begin{exercise}[2] 判断 \(\sum \frac{n!}{n^n}\) 的收敛性。\end{exercise}
\begin{exercise}[2] 判断 \(\sum \frac{1}{\sqrt{n}}\) 的收敛性。\end{exercise}
\begin{exercise}[3] 判断 \(\sum \frac{n^2}{2^n}\) 的收敛性。\end{exercise}
\begin{exercise}[3] 判断 \(\sum \frac{1}{n\ln n}\) 的收敛性。\end{exercise}
\begin{exercise}[3] 判断 \(\sum \frac{\ln n}{n^2}\) 的收敛性。\end{exercise}
\begin{exercise}[3] 判断 \(\sum \frac{n^3}{3^n}\)。\end{exercise}
\begin{exercise}[3] 求 \(\sum_{n=1}^\infty \frac1{n(n+1)}\) 的和（裂项相消）。\end{exercise}
\begin{exercise}[4] 判断 \(\sum \frac{n!}{n^n}\)。\end{exercise}
\begin{exercise}[4] 用根值法判断 \(\sum (\frac{n}{2n+1})^n\)。\end{exercise}
\begin{exercise}[4] 判断 \(\sum \frac{2^n n!}{n^n}\)。\end{exercise}
\begin{exercise}[4] 求 \(\sum \frac1{4n^2-1}\) 的和。\end{exercise}
\begin{exercise}[4] 判断 \(\sum \frac{1}{n^{1+1/n}}\) 的收敛性。\end{exercise}
\begin{exercise}[5] 求 \(\sum \frac{(-1)^{n-1}}{n}\) 的和（交错调和级数）。\end{exercise}
\begin{exercise}[5] 证明 \(\sum \frac{\sin n}{n^2}\) 绝对收敛。\end{exercise}
\begin{exercise}[5] 判断 \(\sum \frac{(-1)^n}{\sqrt{n}}\) 的条件收敛性。\end{exercise}
\begin{exercise}[5] 求 \(\sum \frac1{n^2}\) 的和（提示：\(\pi^2/6\)）。\end{exercise}
\begin{exercise}[6] 证明若 \(\sum a_n\) 绝对收敛则 \(\sum a_n\) 收敛。\end{exercise}
\begin{exercise}[6] 判断 \(\sum \frac{(-1)^n n}{n^2+1}\) 的收敛性。\end{exercise}
\begin{exercise}[6] 证明级数重排定理（黎曼重排定理）。\end{exercise}
\begin{exercise}[7] 证明柯西乘积：\(\sum a_n \cdot \sum b_n = \sum_{n=0}^\infty\sum_{k=0}^n a_k b_{n-k}\)。\end{exercise}
